\documentclass[12pt]{extarticle}
\usepackage[utf8]{inputenc}
\usepackage[margin=1in]{geometry}
\usepackage{cite}
\usepackage{hyperref}
\usepackage{amsmath}
\usepackage{parskip}

\title{\textbf{Beyond Soft PINNs: Differentiable Hard-Constraint Enforcement\\
for Scientific Machine Learning in Julia}}
\author{Angelo Farfan}
\date{March 20, 2026}

\begin{document}
\maketitle

% ─────────────────────────────────────────────
\section*{Motivation and Problem Statement}

Physics-informed neural networks (PINNs) and neural differential equation
models have become a standard way to inject physical knowledge into machine
learning by training neural networks against the residuals of governing
differential equations, boundary conditions, and initial conditions.
Raissi, Perdikaris, and Karniadakis introduced this paradigm and demonstrated
that a single neural network could approximate solutions to complex PDEs when
trained with the right physics-based loss terms~\cite{raissi2019}.
In the Julia ecosystem, this workflow is directly supported by
\texttt{NeuralPDE.jl}~\cite{zubov2021} for PDE problems and by
\texttt{DiffEqFlux.jl}~\cite{rackauckas2019} for neural ordinary differential
equations (neural ODEs), while \texttt{ModelingToolkit.jl}~\cite{rackauckas2022}
provides a symbolic language for specifying governing equations in a clean,
reusable way.

However, a standard PINN is only physics-informed through its \emph{loss
function} --- it does not automatically satisfy physical laws by
construction. Conservation laws, positivity, normalization, or other nonlinear
constraints are typically enforced only through soft penalty terms, and can
therefore remain violated even when the training loss is low. A network may
reduce the residual while still quietly drifting in the quantity you actually
care about, such as total energy or mass. This gap is well-documented in
practice, and has motivated recent work on hard-constraint variants.
Lu~et~al.\ introduced hard-constrained PINNs for inverse design problems,
showing that exact boundary condition enforcement improves both accuracy and
training stability~\cite{lu2021}. More recently, Physics-Constrained Flow
Matching (PCFM) demonstrated that nonlinear projection layers can impose
exact constraints during generative inference without sacrificing model
quality~\cite{utkarsh2025}.

The goal of this project is to build a Julia prototype for
\emph{differentiable hard-constraint enforcement} in scientific machine
learning, with a focus on PINNs and neural ODE models. The central question is:

\begin{quote}
When does replacing a soft constraint penalty with a differentiable
projection or correction step improve physical accuracy, training stability,
and long-horizon prediction quality?
\end{quote}

% ─────────────────────────────────────────────
\section*{Proposed Method}

The proposed approach is to augment a standard PINN or neural ODE pipeline
with a \emph{constraint layer} that maps the network's raw output~$\hat{z}$
to a corrected output~$z^*$ that satisfies one or more physical constraints.
Concretely, the correction solves a small optimization problem at each forward
pass:
\begin{equation}
  z^* = \arg\min_{z}\; \|z - \hat{z}\|^2 \quad \text{subject to} \quad c(z) = 0,
\end{equation}
where $c(z) = 0$ encodes a physical invariant such as energy conservation,
mass conservation, or boundary consistency.
This layer can be placed in one of two settings: (1)~as an \emph{output
correction for PINNs}, projecting the network's predicted solution field onto
the constraint set before evaluating the loss, or (2)~as a
\emph{state correction for neural ODEs}, projecting the integrated state back
to the physically admissible manifold after each time step.

The non-trivial part is that this correction step must be
\emph{differentiable}: gradients need to flow through it during
backpropagation so the network can still be trained end-to-end.
The key insight, formalized in the OptNet work of Amos and
Kolter~\cite{amos2017}, is that instead of differentiating through every
iteration of the inner solver, one can differentiate the optimality
conditions (KKT conditions) of the projection problem directly. By the
implicit function theorem, the gradient of $z^*$ with respect to $\hat{z}$
reduces to solving a single linear system involving the Jacobian of those
optimality conditions --- regardless of how many iterations the forward solver
took. This is both cheaper and more numerically stable than unrolling the
solver.

% ─────────────────────────────────────────────
\section*{Prior Work and Background}

This project sits at the intersection of four lines of existing work.

\textbf{PINNs and NeuralPDE.jl.}
Raissi et al.\ established the PINN training paradigm, showing that PDE
residuals and boundary conditions can be incorporated into the neural network
loss function~\cite{raissi2019}. Zubov et al.\ automated much of this in
Julia with \texttt{NeuralPDE.jl}, supporting symbolic PDE specification,
multiple training strategies, and GPU execution~\cite{zubov2021}.

\textbf{Neural ODEs and DiffEqFlux.jl.}
Chen et al.\ introduced neural ODEs, where an ODE solver is used as a
differentiable layer in the model, making the gradient computation through
the solver a first-class concern~\cite{chen2018}. \texttt{DiffEqFlux.jl}
provides Julia support for these models and can automatically select between
forward-mode, reverse-mode, and adjoint-based differentiation
strategies~\cite{rackauckas2019}.

\textbf{Optimization as a differentiable layer.}
Amos and Kolter demonstrated with OptNet that a constrained quadratic program
can be embedded as a differentiable layer inside a neural network, with the
backward pass derived analytically from the KKT conditions of the
problem~\cite{amos2017}. This is the foundational result that makes the
constraint layer idea computationally tractable.

\textbf{Hard-constraint enforcement.}
Lu et al.\ showed, in the context of inverse design, that hard-constrained
PINNs outperform soft-penalty variants when exact boundary condition
satisfaction is critical~\cite{lu2021}. PCFM extended this idea to
generative models, showing that iterative nonlinear projection can enforce
exact physical constraints during sampling without retraining~\cite{utkarsh2025}.

Together, these works support the view that exact or near-exact constraint
enforcement is a meaningful improvement over soft penalties --- but a general,
reusable Julia interface for doing this inside PINN and neural ODE training
does not yet exist.

% ─────────────────────────────────────────────
\section*{Methodology}

The project will proceed in three stages.

\textbf{Stage 1: Baseline.}
We will implement a standard soft-penalty PINN and neural ODE baseline using
\texttt{NeuralPDE.jl}~\cite{zubov2021} and \texttt{DiffEqFlux.jl}~\cite{rackauckas2019},
with \texttt{ModelingToolkit.jl}~\cite{rackauckas2022} used for symbolic
specification of the governing equations and constraints.
The anchor benchmark will be a Hamiltonian ODE system (e.g., a simple
pendulum or two-body problem) where the physical constraint is energy
conservation. This benchmark is ideal because the failure mode of soft
PINNs --- energy drift over long horizons --- is visually obvious and easy to
quantify.

\textbf{Stage 2: Constraint layer.}
We will implement the correction step using \texttt{Optimization.jl} for the
inner forward solve and \texttt{Zygote.jl} for the backward pass. The first
version will use unrolled differentiation through a fixed number of projection
iterations. If time permits, a second version will use implicit
differentiation through the optimality conditions following the approach of
Amos and Kolter~\cite{amos2017}, requiring a linear solve via
\texttt{LinearSolve.jl} in the backward pass.

\textbf{Stage 3: Comparison.}
We will compare three configurations on the benchmark: (1)~soft PINN
baseline with residual penalties, (2)~projected/corrected model with the
constraint layer, and (3)~a hybrid approach combining soft penalties with a
lighter projection step. The comparison will measure constraint violation,
solution error against a reference solver, training wall time, and training
stability across random seeds.

% ─────────────────────────────────────────────
\section*{Julia Ecosystem Contribution}

The primary software contribution of this project is a reusable module ---
tentatively called \texttt{ConstraintLayers.jl} --- that exposes a small,
clean interface for constrained correction inside scientific ML workflows.
The design philosophy mirrors that of the SciML ecosystem~\cite{rackauckas2022}:
isolate one focused abstraction, connect it to existing Julia packages, and
validate it on concrete examples.
The module would build on \texttt{NeuralPDE.jl}, \texttt{DiffEqFlux.jl},
\texttt{ModelingToolkit.jl}, and \texttt{Optimization.jl} rather than
reimplementing solvers from scratch.

If continued beyond the semester, natural follow-on contributions include: a
generic constraint-layer API that other SciML packages could adopt as an
extension point; documented benchmark examples showing when hard constraints
provide measurable benefit; and a pull request or extension hook into
\texttt{NeuralPDE.jl} or \texttt{DiffEqFlux.jl} once the API stabilizes.

% ─────────────────────────────────────────────
\section*{Evaluation Plan}

The project will evaluate both physical correctness and computational
practicality. The main metrics are: constraint violation after training or
rollout, PDE/ODE residual error, solution error against a reference solver,
training wall time, and training stability across multiple random seeds. This
metric design is consistent with the evaluation approach used in PCFM, which
explicitly measures constraint violation and wall time~\cite{utkarsh2025}, and
with the style of quantitative comparison used in SciML
benchmarks~\cite{rackauckas2022}.

A successful semester-scale outcome is: one solid benchmark problem, one
working constraint layer implementation, and one clear empirical conclusion
about when hard constraints outperform soft penalties. A stretch goal is a
second benchmark or a full implicit backward pass via the KKT conditions.

% ─────────────────────────────────────────────
\section*{Expected Deliverables}

By the end of the semester we expect to produce:
\begin{itemize}
  \item A Julia codebase implementing a first version of the constraint layer,
        with a clean interface over \texttt{Optimization.jl}.
  \item One benchmark study comparing soft vs.\ hard constraint enforcement on
        a Hamiltonian ODE system, with quantitative metrics and training curves.
  \item A short written report covering methodology, results, and discussion.
  \item A roadmap for how this could become a reusable package contribution
        to the SciML ecosystem after the course.
\end{itemize}

% ─────────────────────────────────────────────
\bibliographystyle{plain}
\bibliography{refs}

\end{document}